\begin{proposition}[the integral dictionary]\hypertarget{tag-032R}{}\label{push:src:E.1.1}
\StatementRecord{push:src:E.1.1}{proved}{032R}{\cite[Section 2, equations (2.1), (2.2), (2.7), (2.11)--(2.12)]{ViaclovskyLog2026}}

Let $(\delta,a,b,d)$ be the marking in
\cite[Section 2, equations (2.1), (2.2), (2.7), (2.11)--(2.12)]{ViaclovskyLog2026},
and put $u_0=d+2a+2b$, $u_1=d+3a+3b$.
The homomorphism from the marking of the question to this marking is
\[
 P_V=\begin{pmatrix}
 -2&1&0&1\\-6&0&-1&0\\0&-1&0&0\\-1&0&0&0
 \end{pmatrix},\qquad \det P_V=-1.
\]
It intertwines the two finite monodromies and satisfies
\[
 P_V\widehat\delta=\delta,\qquad
 P_V\varepsilon=-u_0,\qquad
 P_V\varepsilon'=\delta-u_1.
\]
With the counterclockwise deck generators of that paper, its data
$m_0=3k$, $m_1=4l$ and
$a_i=[v_i/m_i]$ with $v_i=\rho_i\delta+N_i u_i$ correspond to
\[
 w_1=N_0\varepsilon-\rho_0\widehat\delta,
 \qquad
 w_2=N_1\varepsilon'-(\rho_1+N_1)\widehat\delta.
\]
For the section-preserving cusp comparison used there, $l_0=0$ and
$p=-(nN_0+mN_1)$.  Thus its sphere condition is
$b:=nN_0+mN_1=\pm1$.
The brackets in $a_i$ mean the torsion point in the smooth
root torus, using its section-preserving origin and the
transported period vector.  Replacing $v_i$ by an integral
multiple $m_i$ times a period does not change that torsion
point, although its chosen lift fixes the collar bookkeeping.
\end{proposition}


\begin{proof}
\begin{enumerate}[label=\arabic*.,leftmargin=*]
\item In Engel's marking the four columns of the change to Viaclovsky's
basis are
\[
 Q_E=\begin{pmatrix}-2&1&0&-2\\-1&0&0&0\\3&-1&1&0\\0&0&0&1\end{pmatrix}.
\]
Direct integral multiplication gives $P_V=Q_E^{-1}P_E$ and
$V_0P_V=P_VA_1$, $V_1P_V=P_VA_2$, where $V_0,V_1$ are the
matrices in that paper's Lemma 2.1.  Multiplication by the three
specified invariant vectors gives the displayed identities.
Inverting a clockwise affine generator replaces its invariant
translation $w_j/m_j$ by $-w_j/m_j$.  Consequently the actual
counterclockwise translation vector is $-P_Vw_j$, which gives the
two formulas, including the otherwise easily missed $N_1$ in the
second coefficient.  Viaclovsky's reference section extends through
the smooth locus of the cusp; its comparison therefore has zero
class in the quotient lattice of toric directions.  This is precisely
$l_0=0$.  Substitution into the definition of $p$ proves its formula.
These matrix identities are independently executed in
\texttt{code/e\_push\_exact.py}, with their output attached.
\end{enumerate}
\end{proof}